\documentclass[
    language=english,
    doctype=recipe,
    institution=none,
]{../../omnilatex}

\RequirePackage{config/document-settings}

\begin{document}

\maketitle

\begin{ingredients}
    \item 2 cups all-purpose flour
    \item 1 teaspoon baking soda
    \item 1/2 teaspoon salt
    \item 1 cup (2 sticks) unsalted butter, softened
    \item 3/4 cup granulated sugar
    \item 3/4 cup packed brown sugar
    \item 1 teaspoon vanilla extract
    \item 2 large eggs
    \item 1 1/2 cups semi-sweet chocolate chips
\end{ingredients}

\begin{instructions}
    \item Preheat the oven to 375\textdegree F (190\textdegree C). Line baking sheets with parchment paper.
    \item In a medium bowl, whisk together the flour, baking soda, and salt. Set aside.
    \item In a large bowl, beat the softened butter, granulated sugar, and brown sugar until light and fluffy, about 2 minutes.
    \item Beat in the vanilla extract and eggs one at a time, mixing well after each addition.
    \item Gradually add the flour mixture to the wet ingredients, mixing on low speed until just combined. Do not overmix.
    \item Fold in the chocolate chips using a spatula.
    \item Drop rounded tablespoonfuls of dough onto the prepared baking sheets, spacing them about 2 inches apart.
    \item Bake for 9--11 minutes, or until the edges are golden brown and the centers are still slightly soft.
    \item Let the cookies cool on the baking sheets for 2 minutes before transferring to wire racks to cool completely.
\end{instructions}

\recipeNote{For the best results, chill the dough for at least 30 minutes before baking. This prevents the cookies from spreading too much and produces a chewier texture. You can also freeze the dough balls for up to 3 months and bake them straight from the freezer, adding 2--3 extra minutes to the bake time.}

\newpage

\recipeMeta{prep}{15 minutes}
\recipeMeta{cook}{20 minutes}
\recipeMeta{servings}{4}
\recipeMeta{difficulty}{Easy}
\recipeMeta{cuisine}{Italian}
\recipeMeta{calories}{420 kcal}

\maketitle

\begin{ingredients}
    \item 400g spaghetti
    \item 4 cloves garlic, thinly sliced
    \item 1/2 teaspoon red pepper flakes
    \item 1/2 cup extra-virgin olive oil
    \item 1 cup fresh parsley, finely chopped
    \item 1 cup freshly grated Parmesan cheese
    \item Salt and black pepper to taste
    \item Juice of 1 lemon (optional)
\end{ingredients}

\begin{instructions}
    \item Bring a large pot of salted water to a rolling boil. Cook the spaghetti according to package directions until al dente. Reserve 1 cup of pasta water before draining.
    \item While the pasta cooks, heat the olive oil in a large skillet over medium heat. Add the sliced garlic and red pepper flakes, saut\'{e}ing until the garlic is golden and fragrant, about 1--2 minutes. Be careful not to burn the garlic.
    \item Add the drained spaghetti to the skillet with the garlic oil. Toss well to coat the pasta evenly.
    \item Add the reserved pasta water a little at a time, tossing continuously, until the sauce is silky and clings to the pasta.
    \item Remove from heat. Stir in the chopped parsley, lemon juice (if using), and half of the Parmesan cheese. Season with salt and pepper to taste.
    \item Serve immediately, topped with the remaining Parmesan cheese and an extra drizzle of olive oil.
\end{instructions}

\recipeNote{The key to great aglio e olio is getting the garlic perfectly golden without burning it. If the garlic starts to brown too quickly, reduce the heat immediately. For a twist, add 1/4 cup of toasted breadcrumbs on top for extra crunch.}

\end{document}
